
\documentclass[11pt]{article}
\usepackage{amsmath,amssymb,amsthm,mathtools,booktabs,enumitem,geometry}
\geometry{margin=1in}
\title{Step 125: Prime-Conductor Finite Source-Frame Model}
\author{RATCHET RH Membrane Program}
\date{May 2026}

\newtheorem{theorem}{Theorem}
\newtheorem{lemma}{Lemma}
\newtheorem{definition}{Definition}
\newtheorem{warning}{Warning}
\newtheorem{corollary}{Corollary}

\begin{document}
\maketitle

\section{Purpose}
This step instantiates the conductor-decoupled source route at the most elementary finite level: complete Dirichlet character families modulo a prime conductor.  The goal is not to re-derive approximate functional equations, mollifier length bounds, or large-sieve technology.  Those are imported analytic-number-theory records.  The goal is to translate the finite orthogonality theorem into the membrane lower-frame language, then expose the defects that remain before any completed RH claim can be made.

\section{Setup}
Let $q$ be prime and let $\mathcal N\subset\{1,\ldots,L\}$ with $L<q$.  Let $d=|\mathcal N|$.  For each Dirichlet character $\chi\bmod q$ define
\[
Q_\chi a=\sum_{n\in\mathcal N}a_n\chi(n).
\]
The finite source Gram is
\[
G_q(m,n)=\sum_{\chi\bmod q}\chi(n)\overline{\chi(m)}.
\]
Because $q$ is prime and $L<q$, every $n\in\mathcal N$ is a unit and $m\equiv n\pmod q$ if and only if $m=n$.

\section{Full-character frame}
\begin{theorem}[Prime full-character finite frame]
For $q$ prime and $L<q$,
\[
\sum_{\chi\bmod q} Q_\chi^*Q_\chi=(q-1)I_{\mathbb C^{\mathcal N}}.
\]
Equivalently,
\[
G_q=(q-1)I.
\]
\end{theorem}

\begin{proof}
This is the standard orthogonality relation for the complete character group of $(\mathbb Z/q\mathbb Z)^\times$:
\[
\sum_{\chi\bmod q}\chi(n)\overline{\chi(m)}=
\begin{cases}
q-1,& n\equiv m\pmod q,\\
0,& n\not\equiv m\pmod q.
\end{cases}
\]
Since $1\le m,n<q$, congruence is equality.
\end{proof}

Thus the clean finite coefficient-frame prototype has lower-frame strength
\[
\gamma_q^{\rm full}=q-1
\]
if one uses the raw cumulative source sum.  If the family is normalized by averaging, the strength is $1$ instead.

\section{Primitive-character removal}
For prime $q$, the primitive characters are exactly the nonprincipal characters.  Removing the principal character subtracts the rank-one all-ones matrix
\[
J_d=\mathbf 1\mathbf 1^*.
\]
Thus
\[
G_q^{\rm prim}=(q-1)I-J_d.
\]
Its eigenvalues are
\[
q-1\quad(d-1\text{ times}),
\qquad
q-1-d\quad(1\text{ time}).
\]
Therefore
\[
G_q^{\rm prim}\succeq(q-1-d)I.
\]
This remains positive exactly when $d<q-1$.  Under average normalization by $q-2$, the lower-frame constant is
\[
\frac{q-1-d}{q-2}=1-O(d/q).
\]
So primitive removal is harmless only if the coefficient dimension is $o(q)$.

\section{The source-normalization fork}
The membrane lower-frame theorem uses
\[
F_N=\sum_{\omega\in\mathcal X_N}\lambda_{\omega,N}Q_\omega^*\Theta_\omega^{-1}Q_\omega.
\]
For prime characters, choosing $\lambda_{\chi,N}=1$ gives a raw cumulative source frame with strength $\gamma_N\asymp q_N$.  Choosing $\lambda_{\chi,N}=1/\varphi(q_N)$ gives an averaged frame with strength $O(1)$.

\begin{warning}[Normalization is load-bearing]
The divergence $\Lambda_N\to\infty$ is not a consequence of finite character orthogonality alone.  It depends on the audit currency assigned to source records.  Raw cumulative sources can give $\gamma_N\to\infty$; averaged families do not.  This must be declared as a $P_6$ audit record, not chosen after the proof needs growth.
\end{warning}

\section{Conductor-decoupled finite prototype}
Let the regularized Burnol--Dirichlet shadow have effective coefficient length
\[
L_N\asymp X_NR_N.
\]
Choose a prime conductor
\[
q_N>L_N.
\]
Then the finite full-character frame gives
\[
G_{q_N}\succeq(q_N-1)I.
\]
The primitive frame gives
\[
G_{q_N}^{\rm prim}\succeq(q_N-1-L_N)I.
\]
If $L_N=q_N^\sigma$ with $\sigma<1$, then the primitive averaged lower-frame constant tends to $1$, while the raw primitive lower-frame strength grows like $q_N$.

\section{Interface with the residual visibility gate}
From Step 123,
\[
c_{\rm hyb,N}\ge 1-E_N^2,
\]
where
\[
E_N=\epsilon_{B,N}+\epsilon_{\rm Mell,N}+\tau_{M,N}+\kappa_{\rm tail,N}.
\]
Thus finite residual source absorption follows from
\[
R_N^*G_{q_N}R_N\succeq \gamma_N c_{\rm hyb,N}G_{R,N}.
\]
A positive floor is sufficient:
\[
E_N\le e<1,
\qquad
\gamma_N\to\infty
\quad\Longrightarrow\quad
\gamma_Nc_{\rm hyb,N}\to\infty.
\]

\section{Honest defects}
The prime-conductor model has the following first defects:
\begin{enumerate}[leftmargin=2em]
\item \textbf{Aliasing defect.} If $L_N\ge q_N$, then distinct coefficients can be congruent modulo $q_N$, producing off-diagonal collisions.
\item \textbf{Primitive removal defect.} Passing from all characters to primitive characters subtracts $J_d$ and costs the one-dimensional all-ones direction.
\item \textbf{Composite-modulus oldform defect.} Prime conductors avoid imprimitive conductor stratification. Composite moduli require a separate oldform/imprimitive audit.
\item \textbf{Source-weight normalization defect.} Raw and averaged source currencies give different growth behavior.
\item \textbf{Finite-to-completed tail.} The finite coefficient frame is RH-relevant only after the residual ledger is fixed or exhaustive.
\item \textbf{Matrix moment defect.} If the source frame is $L$-weighted, mollified, or restricted, scalar lower moments are insufficient; one needs a matrix lower-frame theorem.
\end{enumerate}

\section{What is imported versus new}
The finite orthogonality relation is standard.  Approximate functional equations, smooth Muentz regularization, mollifier-length restrictions, and conductor-aspect bilinear forms are imported analytic-number-theory inputs.  The project-native contribution is the membrane reading:
\[
\text{finite source frame} + \Xi\text{-visibility} + \text{fixed/exhaustive tail} \Rightarrow \text{residual absorption}.
\]

\section{Next target}
The next step is not another re-derivation of character orthogonality.  It is the source-weighted matrix lower-frame theorem:
\[
G_{\mathcal X,N}^{\rm weighted}\succeq\gamma_NH_N
\]
for the exact coefficient windows produced by the regularized Burnol--Dirichlet shadow.  Existing large-sieve and hybrid moment results should be treated as imported technology and audited for whether they provide lower-frame, not merely upper-bound or scalar-moment, content.

\end{document}
